\providecommand{\Nb}{\mathcal{N}}
\providecommand{\GL}{\operatorname{GL}}
\providecommand{\Aut}{\operatorname{Aut}}

\chapter{Geometric Priors: Invariance and Equivariance}
\label{ch:priors}

The previous chapters introduced the mathematical language of geometric deep
learning (GDL) --- groups, representations, and the geometric domains on which
data live --- and argued that the failure of classical architectures on
non-Euclidean data is conceptual rather than technical. This chapter turns that
diagnosis into a constructive principle. The central thesis is simple to state:
a model should respect the symmetries of the problem it solves, and the way to
guarantee this is to build the symmetry into the architecture rather than to
hope it will be learned from data. The two notions that make this precise are
\emph{invariance} and \emph{equivariance}. From them we derive a design
blueprint --- the \emph{geometric priors} of the title --- that unifies
convolutional, attention-based, and graph architectures under a single
geometric construction.

\section{Why traditional networks need a geometric prior}
\label{sec:priors:motivation}

A traditional artificial neural network treats its input as an unstructured
vector $x \in \R^n$. A single neuron computes $y = \sigma(\inner{w}{x} + b)$,
combining all $n$ coordinates through a global weighted sum. Reinterpreted as an
operator on a signal over a domain $\Omega = \{1, \dots, n\}$, this domain is a
bare index set: it carries no notion of \emph{locality} (every coordinate
contributes equally, with no concept of near and far), no \emph{structured
neighbourhood} (the connection $w_1 x_1$ is structurally identical to
$w_{100} x_{100}$), and no \emph{geometric invariance} (permuting the indices
generally changes the output, even when the permutation is semantically
irrelevant). The network ignores any geometry implicit in the data.

This blindness has a concrete cost. Suppose the quantity we wish to predict is
unchanged under a group $\group$ of transformations --- rotations of an image,
relabellings of the nodes of a graph. A fully connected network has no built-in
reason to produce the same answer for $x$ and for a transformed input
$g \cdot x$. The standard remedy, \emph{data augmentation}, presents the network
with transformed copies of every example; but this inflates the training set by
a factor of $\abs{\group}$, offers no guarantee that the learned function is
actually invariant, and degrades or fails entirely when $\group$ is large or
continuous. A geometric prior replaces this hope with a guarantee: the symmetry
is encoded in the structure of the map, so it holds exactly, for every input,
by construction.

\section{Invariant and equivariant functions}
\label{sec:priors:inv-equi}

We work with a group $\group$ acting on the relevant spaces. Recall that an
\emph{action} of $\group$ on a set $X$ is a map $\group \times X \to X$,
written $(g, x) \mapsto g \cdot x$, that is compatible with the group law:
$e \cdot x = x$ and $g \cdot (h \cdot x) = (gh) \cdot x$. When $X$ is a vector
space and each $g$ acts linearly, the action is a \emph{representation}
$\rho_X \colon \group \to \GL(X)$ (\Cref{ch:foundations}). We write $g \cdot x$
for $\rho_X(g)\, x$ throughout.

\begin{definition}[Equivariant function]
\label{def:priors:equivariance}
Let $\group$ be a group acting on $X$ via $\rho_X$ and on $Y$ via $\rho_Y$.
A map $f \colon X \to Y$ is \emph{$\group$-equivariant} if
\[
  f(g \cdot x) = g \cdot f(x)
  \qquad \text{for all } g \in \group,\ x \in X,
\]
that is, $f \circ \rho_X(g) = \rho_Y(g) \circ f$ for every $g$.
\end{definition}

\begin{definition}[Invariant function]
\label{def:priors:invariance}
A map $f \colon X \to Y$ is \emph{$\group$-invariant} if
$f(g \cdot x) = f(x)$ for all $g \in \group$ and $x \in X$.
\end{definition}

Invariance is the special case of equivariance in which $\group$ acts trivially
on the output, $\rho_Y(g) = \mathrm{id}$. Equivariance is the more fundamental
notion: it says that transforming the input produces a correspondingly
transformed output, so that geometric structure is \emph{preserved} rather than
discarded. Equivariance is the property we want from the intermediate layers of
a network, where structure must survive; invariance is appropriate only at the
final readout, where we collapse to a label.

\Cref{fig:priors:equivariance} renders \Cref{def:priors:equivariance} as a
commutative square: equivariance is exactly the statement that the two paths
from $X$ to $Y$ --- transform then apply $f$, or apply $f$ then transform ---
agree.

\begin{figure}[h]
\centering
\begin{tikzpicture}[
    node distance=2.6cm and 4cm,
    space/.style={circle, draw=gdlgray!80, fill=gdlblue!15, minimum size=1.5cm, font=\large},
    arrow/.style={-{Stealth[length=7pt]}, thick},
    lbl/.style={font=\small, midway}
]
\node[space] (X) {$X$};
\node[space, right=of X] (Y) {$Y$};
\node[space, below=of X] (gX) {$X$};
\node[space, below=of Y] (gY) {$Y$};
\draw[arrow, gdlblue!70!black] (X) -- node[lbl, above] {$f$} (Y);
\draw[arrow, gdlblue!70!black] (gX) -- node[lbl, below] {$f$} (gY);
\draw[arrow, gdlred!70!black] (X) -- node[left, align=center] {\small\itshape transform\\[1pt] $\rho_X(g)$} (gX);
\draw[arrow, gdlred!70!black] (Y) -- node[lbl, right] {$\rho_Y(g)$} (gY);
\node[font=\Large] at ($(X)!0.5!(gY)$) {$\circlearrowleft$};
\node[font=\normalsize, below=0.4cm of gX, anchor=north west, xshift=0.6cm] {
    $f(\rho_X(g)\cdot x) = \rho_Y(g)\cdot f(x)$
};
\end{tikzpicture}
\caption[The equivariance commutative square]{Equivariance as a commutative
diagram. A map $f \colon X \to Y$ is $\group$-equivariant when the square
commutes: transforming the input by $\rho_X(g)$ and then applying $f$ yields the
same result as applying $f$ and then transforming the output by $\rho_Y(g)$.}
\label{fig:priors:equivariance}
\end{figure}

\begin{example}[Three canonical symmetries]
\label{ex:priors:symmetries}
Image classification is invariant to translations: shifting the picture of a cat
does not change the label, so the score function should be
$\Z^2$-invariant. Semantic segmentation, by contrast, is
\emph{equivariant} to translations: shifting the image shifts the segmentation
mask by the same amount. A function defined on the nodes of a graph should be
equivariant to the permutation group $S_n$: relabelling the nodes permutes the
outputs identically, because a graph carries no canonical ordering of its
vertices. Each task fixes a group $\group$ and a required behaviour --- invariant
or equivariant --- under it.
\end{example}

The reason equivariance is the right primitive for deep architectures is that it
is preserved by composition.

\begin{theorem}[Composition of equivariant maps]
\label{thm:priors:composition}
Let $\group$ act on $X$, $Y$, $Z$ via $\rho_X, \rho_Y, \rho_Z$. If
$f \colon X \to Y$ and $h \colon Y \to Z$ are $\group$-equivariant, then so is
the composition $h \circ f \colon X \to Z$.
\end{theorem}

\begin{proof}
Fix $g \in \group$ and $x \in X$. Using equivariance of $f$ and then of $h$,
\[
  (h \circ f)(g \cdot x) = h\bigl(f(g \cdot x)\bigr)
    = h\bigl(g \cdot f(x)\bigr)
    = g \cdot h\bigl(f(x)\bigr)
    = g \cdot (h \circ f)(x). \qedhere
\]
\end{proof}

\Cref{thm:priors:composition} is the keystone of the whole framework. It
guarantees that a network built by stacking equivariant layers is itself
equivariant, end to end --- equivariance is a property we establish once, per
layer, and then inherit for free at every depth. Composing such an equivariant
backbone with a single invariant readout at the end yields a globally invariant
predictor. This turns symmetry from a property to be checked \emph{a posteriori}
into a constructive recipe for entire architectures.

\section{The geometric deep learning blueprint}
\label{sec:priors:blueprint}

The composition principle suggests a template for designing architectures on any
geometric domain~\citep{bronstein2021geometric}. Fix a domain $\Omega$ (a grid,
a graph, a manifold), a space of signals $\mathcal{X}(\Omega)$ on it, and a
symmetry group $\group$ acting on both. The \emph{geometric deep learning
blueprint} assembles a model from four kinds of building block.

\begin{description}
\item[Equivariant layers.] Linear or nonlinear maps
$\mathcal{X}(\Omega) \to \mathcal{X}(\Omega')$ that commute with the action of
$\group$. These are the workhorses: they extract features while transforming
predictably under the symmetry. A pointwise nonlinearity applied after an
equivariant linear map preserves equivariance as long as it acts identically on
every position.
\item[Local pooling (coarsening).] Equivariant maps that reduce the resolution
of the domain, $\Omega \to \Omega'$, while respecting the symmetry. Pooling
implements \emph{scale separation}: it lets later layers see larger
neighbourhoods through a coarser domain, building a hierarchy of features.
\item[Invariant global readout.] A final $\group$-invariant map that collapses
the signal to the desired output --- a single label, a graph-level property ---
typically a symmetric aggregation (sum, mean, max) followed by a standard
network.
\item[Symmetry-respecting attention or message functions.] When connectivity
itself should adapt to content, the weights linking positions are computed by
functions that are themselves equivariant. This generalises fixed local filters
to data-dependent ones.
\end{description}

Two structural principles underlie this template and recur throughout the book:

\begin{description}
\item[Symmetry (equivariance).] Every internal operation commutes with
$\group$, so the global function is equivariant by
\Cref{thm:priors:composition}, with invariance achieved only at the readout.
\item[Scale separation (locality and compositionality).] Operations are
\emph{local} with respect to the geometry of $\Omega$ --- each layer aggregates
only over a neighbourhood --- and complex representations are built
\emph{compositionally}, by stacking local layers interleaved with coarsening.
Locality keeps the parameter count and computation tractable; compositionality
lets a deep stack of local operations capture long-range, global structure.
\end{description}

\begin{remark}[One blueprint, five domains]
\label{rem:priors:fivegs}
The blueprint is deliberately domain-agnostic. Instantiating it on a grid with
the translation group recovers convolutional networks (\Cref{ch:grids});
extending the group to rotations and reflections gives group-equivariant
networks (\Cref{ch:groups}); replacing the grid by an arbitrary graph with the
permutation group gives graph neural networks (\Cref{ch:graphs}); endowing the
domain with intrinsic distances gives learning on manifolds
(\Cref{ch:geodesics}); and acknowledging the absence of a canonical local frame
on a curved domain gives gauge-equivariant networks (\Cref{ch:gauges}). The
``five Gs'' are five concrete realisations of one geometric idea.
\end{remark}

\section{Universality and expressivity}
\label{sec:priors:universality}

A natural worry is that constraining a network to be equivariant might cripple
its expressive power: the constraint carves out a strict subspace of all
functions. The following result, which extends the classical universal
approximation theorem to the equivariant setting, dispels that worry. It is due
to \citet{yarotsky2022} and \citet{kondor2018generalizationconvolutionalneural}.

\begin{theorem}[Equivariant universal approximation]
\label{thm:priors:universal}
Let $\group$ be a compact group acting continuously on compact metric spaces $X$
and $Y$, and let $\mathcal{F}^{\group}(X, Y)$ be the space of all continuous
$\group$-equivariant maps $f \colon X \to Y$. Then for every
$f \in \mathcal{F}^{\group}(X, Y)$ and every $\varepsilon > 0$ there exists a
$\group$-equivariant neural network $\mathcal{N}$ such that
\[
  \sup_{x \in X} d_Y\bigl(f(x), \mathcal{N}(x)\bigr) < \varepsilon,
\]
where $d_Y$ is the metric on $Y$.
\end{theorem}

\begin{proof}[Proof sketch]
The argument has two ideas. First, by the classical universal approximation
theorem there is an unconstrained network $\tilde f$ with
$\sup_x d_Y(f(x), \tilde f(x)) < \varepsilon$. Second, we symmetrize $\tilde f$
through the \emph{Reynolds operator}, averaging over the group with respect to
the normalized Haar measure $\mu$:
\[
  \bar f(x) = \int_{\group} \rho_Y(g^{-1})\, \tilde f(g \cdot x)\, d\mu(g).
\]
A change of variables together with the invariance of the Haar measure shows
$\bar f(h \cdot x) = \rho_Y(h)\, \bar f(x)$, so $\bar f$ is equivariant. Finally,
since $f$ is already equivariant and (for a compact group) the representations
may be taken unitary, $\rho_Y(g^{-1})$ preserves distances and the averaging
does not increase the error:
$\sup_x d_Y(\bar f(x), f(x)) \le \sup_x d_Y(\tilde f(x), f(x)) < \varepsilon$.
\end{proof}

The message is decisive: \emph{equivariance costs nothing in expressivity}. A
$\group$-equivariant network can approximate any continuous $\group$-equivariant
target to arbitrary accuracy. The symmetry constraint does not shrink the set of
reachable functions within the relevant subspace; it acts purely as an inductive
bias, discarding the irrelevant non-equivariant functions and so improving
sample efficiency and generalization without sacrificing capacity. The empirical
success of convolutional networks --- equivariant to the translation group
$\Z^2$ --- is a corollary~\citep{cohen2016group}.

\section{Message passing: the unifying framework}
\label{sec:priors:message-passing}

The blueprint becomes algorithmically concrete through a single computational
primitive: \emph{message passing}. Viewing data as a graph
$\mathcal{G} = (V, E)$ of nodes and edges, a layer propagates information by
having each node aggregate transformed features from its neighbours. This one
construction subsumes the major deep-learning architectures; their apparent
differences reduce to choices of graph topology, connectivity, and symmetry
group.

\begin{definition}[Message-passing neural network]
\label{def:priors:mpnn}
A \emph{message-passing neural network} (MPNN) operates on a graph
$\mathcal{G} = (V, E)$ with node features $\{h_v\}_{v \in V}$ and edge features
$\{e_{uv}\}$. Layer $\ell$ updates each node in two steps:
\[
  m_v^{(\ell+1)} = \bigoplus_{u \in \Nb(v)}
      M_\ell\bigl(h_v^{(\ell)}, h_u^{(\ell)}, e_{vu}\bigr),
  \qquad
  h_v^{(\ell+1)} = U_\ell\bigl(h_v^{(\ell)}, m_v^{(\ell+1)}\bigr),
\]
where $\Nb(v)$ is the neighbourhood of $v$, $M_\ell$ a learnable message
function, $\bigoplus$ a permutation-invariant aggregator (sum, mean, max), and
$U_\ell$ an update function.
\end{definition}

The aggregator $\bigoplus$ is permutation-invariant by design, which makes each
layer equivariant to relabellings of the neighbours --- the message-passing
layer is the blueprint's equivariant layer for relational data. The following
correspondences (\Cref{fig:priors:mpnn}) show that the familiar architectures
are special cases.

\begin{proposition}[Architectures as message passing]
\label{prop:priors:architectures}
The following are instances of \Cref{def:priors:mpnn}:
\begin{enumerate}
\item A fully connected layer $h^{(\ell+1)} = \sigma(W h^{(\ell)} + b)$ is an
MPNN on the complete graph $K_d$ whose features are nodes, with a pair-specific
message $M(h_v, h_u) = W_{vu} h_u$ and sum aggregation. There is no weight
sharing and hence no geometric prior.
\item A convolution on a regular grid $\Z^d$ with kernel support $\mathcal{K}$
is an MPNN with neighbourhood $\Nb(v) = \{u : v - u \in \mathcal{K}\}$ and
message $M(h_v, h_u) = W_{v-u}\, h_u$: the weight depends only on the
\emph{relative} position $v - u$, not on the absolute pair $(v, u)$. This weight
sharing is exactly translation equivariance, and it reduces the parameter count
from $O(n^2 c^2)$ to $O(\abs{\mathcal{K}}\, c^2)$.
\item A self-attention layer on $n$ tokens is an MPNN on the complete graph
$K_n$ with message $M(h_i, h_j) = \alpha_{ij}\, V h_j$, where the attention
weights $\alpha_{ij} = \mathrm{softmax}_j\bigl((Q h_i)^{\!\top} (K h_j) /
\sqrt{d_k}\bigr)$ are computed from the node \emph{content}, not from position,
and aggregation is the convex combination $m_i = \sum_j \alpha_{ij} V h_j$.
\item A graph neural network is the general case: an MPNN on an arbitrary graph
$\mathcal{G}$, equivariant to its automorphism group.
\end{enumerate}
\end{proposition}

Reading across these cases reveals a single axis of increasing geometric
sophistication: fully connected networks impose no prior (uniform, pair-specific
weights); convolutions impose a \emph{spatial} prior (weights shared by relative
position, encoding translation equivariance); transformers replace fixed weights
by \emph{adaptive, content-dependent} ones (global connectivity at quadratic
cost); and graph networks generalise the whole scheme to arbitrary relational
topologies. The geometric prior is precisely what each architecture assumes
about which connections matter and how their weights are shared.

\begin{figure}[h]
\centering
\begin{tikzpicture}[>=Stealth, font=\small,
    nd/.style={circle, draw, thick, fill=white, minimum size=4mm, inner sep=0pt},
    edge/.style={thick, draw=gdlgray!60},
    box/.style={rectangle, rounded corners=4pt, draw=gdlgray!70, thick, minimum width=4.6cm, minimum height=3.7cm},
    title/.style={font=\small\bfseries, fill=white, inner sep=2pt},
    info/.style={font=\scriptsize, text=black!70, align=center}
]
% ANN
\begin{scope}[shift={(-3.3, 2.7)}]
    \node[box, fill=gdlblue!15] (annbox) {};
    \node[title] at (0, 2.05) {ANN (fully connected)};
    \foreach \i/\a in {1/90,2/162,3/234,4/306,5/18} \node[nd, fill=gdlblue!20] (a\i) at (\a:0.9) {};
    \foreach \i in {1,...,5} \foreach \j in {1,...,5} \ifnum\i<\j \draw[edge, gdlblue!30] (a\i)--(a\j); \fi
    \node[info] at (0, -1.45) {complete graph $K_d$\\specific weights $W_{ij}$};
\end{scope}
% CNN
\begin{scope}[shift={(3.3, 2.7)}]
    \node[box, fill=gdlgreen!15] (cnnbox) {};
    \node[title] at (0, 2.05) {CNN (convolution)};
    \foreach \i in {0,1,2} \foreach \j in {0,1,2} \node[nd, fill=gdlgreen!20] (c\i\j) at ({\i*0.7-0.7},{\j*0.7-0.7}) {};
    \foreach \i in {0,1,2} \foreach \j [evaluate=\j as \jn using int(\j+1)] in {0,1} {
        \draw[edge, gdlgreen!50] (c\i\j)--(c\i\jn); \draw[edge, gdlgreen!50] (c\j\i)--(c\jn\i);}
    \draw[thick, gdlorange, rounded corners=2pt] (-0.3,-0.3) rectangle (0.3,0.3);
    \node[info] at (0, -1.45) {grid $\Z^d$, shared weights\\$\group=(\Z^d,+)$};
\end{scope}
% Transformer
\begin{scope}[shift={(-3.3, -2.7)}]
    \node[box, fill=gdlpurple!15] (transbox) {};
    \node[title] at (0, 2.05) {Transformer (attention)};
    \foreach \i/\a in {1/90,2/162,3/234,4/306,5/18} \node[nd, fill=gdlpurple!20] (t\i) at (\a:0.9) {};
    \draw[edge, gdlpurple!60, line width=2pt] (t1)--(t3);
    \draw[edge, gdlpurple!40, line width=1.4pt] (t1)--(t2);
    \draw[edge, gdlpurple!20, line width=0.5pt] (t1)--(t4);
    \draw[edge, gdlpurple!50, line width=1.4pt] (t2)--(t4);
    \draw[edge, gdlpurple!60, line width=2pt] (t3)--(t5);
    \draw[edge, gdlpurple!30, line width=0.8pt] (t4)--(t5);
    \node[info] at (0, -1.45) {complete graph $K_n$\\adaptive weights $\alpha_{ij}$};
\end{scope}
% GNN
\begin{scope}[shift={(3.3, -2.7)}]
    \node[box, fill=gdlred!15] (gnnbox) {};
    \node[title] at (0, 2.05) {GNN (graph network)};
    \node[nd, fill=gdlred!20] (g1) at (-0.7,0.6) {};
    \node[nd, fill=gdlred!20] (g2) at (0.3,0.8) {};
    \node[nd, fill=gdlred!20] (g3) at (0.8,0.1) {};
    \node[nd, fill=gdlred!20] (g4) at (-0.4,-0.3) {};
    \node[nd, fill=gdlred!20] (g5) at (0.4,-0.5) {};
    \node[nd, fill=gdlred!20] (g6) at (-0.15,0.25) {};
    \foreach \a/\b in {g1/g2,g1/g4,g1/g6,g2/g3,g2/g6,g3/g5,g4/g5,g4/g6,g5/g6}
        \draw[edge, gdlred!50] (\a)--(\b);
    \node[info] at (0, -1.45) {arbitrary graph $\mathcal{G}$\\$\group=\Aut(\mathcal{G})$};
\end{scope}
% Centre
\node[font=\bfseries, align=center, fill=gdlyellow!20, rounded corners=3pt,
      draw=gdlyellow!70, thick, inner sep=5pt] (mpnn) at (0,0)
      {MPNN\\[-1pt]{\scriptsize $m_v=\bigoplus_{u\in\Nb(v)} M(h_v,h_u,e_{vu})$}};
\draw[->, thick, gdlblue!60] (annbox.south east) -- (mpnn.north west);
\draw[->, thick, gdlgreen!50!black] (cnnbox.south west) -- (mpnn.north east);
\draw[->, thick, gdlpurple!60] (transbox.north east) -- (mpnn.south west);
\draw[->, thick, gdlred!60] (gnnbox.north west) -- (mpnn.south east);
\end{tikzpicture}
\caption[Architectures as message passing]{The major deep-learning
architectures as instances of message passing. They differ only in graph
topology, weight-sharing scheme, and symmetry group $\group$: fully connected
networks (complete graph, pair-specific weights), convolutional networks (grid,
translation-shared weights), transformers (complete graph, content-adaptive
weights), and graph networks (arbitrary graph). The message-passing update at
the centre is common to all.}
\label{fig:priors:mpnn}
\end{figure}

\subsection{Expressive power and its limits}
\label{sec:priors:expressivity-limits}

Message passing is unifying but not omnipotent. Its expressive power on graphs is
pinned down by a classical combinatorial algorithm, and its depth is limited by
two complementary pathologies. We sketch these here at a high level; the
spectral and geometric analysis returns in \Cref{ch:graphs,ch:geodesics}.

\paragraph{The Weisfeiler--Leman ceiling.}
The Weisfeiler--Leman test (1-WL) is a graph-isomorphism heuristic that
iteratively refines a colouring of the nodes: each node's new colour hashes its
old colour together with the multiset of its neighbours'
colours~\citep{weisfeiler1968reduction}. \citet{xu2019how} and
\citet{morris2019weisfeiler} proved that an MPNN with injective aggregation and
update functions is \emph{exactly} as discriminative as 1-WL: it can never
distinguish two graphs that 1-WL deems equivalent, and a suitable MPNN (the
Graph Isomorphism Network) attains this bound. Consequently MPNNs cannot
distinguish certain non-isomorphic graphs --- for instance a six-cycle $C_6$
from two disjoint triangles $K_3 \sqcup K_3$, which share the same degree
sequence and local colour histogram --- exposing a fundamental ceiling on what
local message passing can perceive. Operating on higher-order structures
(simplicial complexes) or adding structural and positional encodings can break
this barrier.

\paragraph{Oversmoothing and oversquashing.}
Stacking many message-passing layers degrades performance in two opposite ways.
\emph{Oversmoothing} is the loss of discriminative signal: a linear MPNN acts as
discrete heat diffusion $H^{(\ell+1)} = (I - \alpha \tilde L)\, H^{(\ell)}$
driven by the graph Laplacian $\tilde L$, and the Dirichlet energy of the node
features decays exponentially, $E(H^{(\ell)}) \le (1 - \alpha
\lambda_2)^{2\ell}\, E(H^{(0)})$, so representations converge to a constant and
nodes become indistinguishable~\citep{li2018deeper}; the rate is controlled by
the spectral gap $\lambda_2$. \emph{Oversquashing} is the dual failure: long-range
information must be compressed through topological bottlenecks, and the
sensitivity $\|\partial h_v^{(\ell)} / \partial h_u^{(0)}\|$ of a node to a
distant one decays with the powers of the transition matrix, vanishing entirely
beyond the graph distance $\ell$. \citet{topping2022understanding} relate these
bottlenecks to negative discrete Ricci curvature, a striking instance of
differential geometry diagnosing --- and prescribing cures for --- a limitation
of neural networks.

These limits do not undermine the framework; they sharpen it. They are
manifestations of the fundamentally \emph{local} character of message passing,
and the geometric tools developed in the rest of this book --- spectral
analysis, discrete curvature, higher-order topology --- both explain them and
suggest remedies.

\section{Summary}
\label{sec:priors:summary}

Invariance and equivariance turn the informal demand that a model ``respect the
symmetries of the data'' into precise constraints on a function. Because
equivariant maps compose (\Cref{thm:priors:composition}), these constraints
yield a constructive blueprint: stack equivariant layers and local coarsening,
finish with an invariant readout. The equivariant universal approximation
theorem (\Cref{thm:priors:universal}) certifies that this discipline costs no
expressive power. Message passing makes the blueprint concrete and reveals
convolutional, attention-based, and graph architectures as one construction
under different geometric priors. The remaining chapters instantiate this
blueprint on each of the five geometric domains.
